\section{Protocol history}
\label{app:history}

Every change below is recorded, with its date, in the protocol documents released with the code. None was
made after looking at confirmation data.

\paragraph{v1 (frozen design, never run).} The design fixed the hypotheses, estimands, factorial,
constructions, gates, bootstrap and held-out confirmation template. An audit before any inference found
two problems. (i) The action events began with a space (\lit{\textvisiblespace APPROVE}). Qwen3-8B begins
its turn without one, and the two strings tokenize differently, so v1 would have scored continuations the
model never produces. (ii) Among the H3 predictors was the decision's own action confidence, which is not
available before the decision is made.

\paragraph{v1b (corrected before inference).} Action events became the label exactly as generated plus the
end-of-turn token. The H3 predictors were split into pre-decision predictors and a contemporaneous
reference. Historical retrieval was added as a competence measure. During setup, the runtime was changed
twice for numerical reasons (Appendix~\ref{app:precision}), and a rounding tolerance of $10^{-5}$ was added
to validators that rejected float32 log probabilities a few $10^{-7}$ above zero. The v1b gate failed:
direct retrieval was 0.851--0.946 in every cell. The cause was the prompt format. The user message ended
with a literal \lit{Answer:} line, which the model often echoed (about 260 of 4{,}224 retrieval answers);
about 100 answers gave an action label instead of a state. Only 64 named the obsolete value. The decision
effects replicated on the gate histories.

\paragraph{v2 (first redesign; at most two were allowed).} The redesign changed only what the failure
required: it removed the trailing \lit{Answer:} line and used new seeds and entity-name pools so that all
histories were new. Confirmation size was set to four replicates (768 histories) before any v2 data
existed, to fit the compute budget. The scorer was changed from a batched scorer to a key--value-cache
scorer for speed, before any v2 effect was inspected. Both scorers compute the same quantity, and the
cache scorer passed its equivalence check. The v2 gate passed in three cells, and confirmation data were
then generated.

\paragraph{Second model.} The model list, adding Gemma 3 4B with the same rows, rules and runtime, was
fixed before any Gemma data were scored. A 24-record smoke test, inspected for format only, showed that
Gemma writes its label, a newline and then its end-of-turn token. The scored events had omitted the
newline and covered at most a few percent of the probability mass. The scored event was corrected to include the
newline, the checks were rerun, and the smoke outputs were set aside. Development and gate then failed the
competence criteria in every cell (\S\ref{sec:results-gemma}).

\paragraph{Replication models.} A further amendment, dated before any of their data were generated, listed
five models in a fixed order with pinned revisions. It fixed the per-model sequence, the competence
criteria, the reporting rule, and the definition of competent cells: those passing on the model's own
frozen gate. For this submission the scope was reduced, before any of the five had produced data, to the
first model, Qwen3-14B. The other four are deferred under the same amendment. Qwen3-14B's preflight needed
no correction (causality 0.0; scorer difference $8.3\times10^{-7}$ nats; it ends its answer with the
label and its end-of-turn token). Its rows were byte-identical to Qwen3-8B's in every split.

\subsection{Numerical causality of log-probability scoring}
\label{app:precision}

Scoring answers by teacher forcing assumes that the log probability at a position does not depend on
tokens appended after it. We tested this directly on the same prompts, by comparing the next-token
distribution at a fixed position with and without one or two tokens appended. With 8-bit weights
(LLM.int8(), \citealp{dettmers2022llmint8}), log probabilities changed by up to 0.843 nats. With bfloat16, the
top-20 log probabilities changed by up to 0.875 nats, and computing only the output layer in float32 reduced it to
0.502. In float32, and with 4-bit nf4 weights, it was exactly zero. In practice this broke our
validators: four disjoint candidate answers summed to more than probability one. A key--value-cache
scorer also disagreed with full recomputation by 3.406 nats under 8-bit weights, but by
$5.08\times10^{-5}$ in float32. Score differences of a few nats are the size of the effects studied here,
so we ran all inference in float32. Score-based studies that use reduced precision should report a
check of this kind.
